% New abstract theorem. Not a geometric nonalgebraicity assertion.
% Certified abstract theorem; geometric transfer remains open.
\paragraph{C.4z (Proved): a correlated graph obstruction at every level.}
Put $\Lambda=\mathbb Z_{(2)}$, $U=\Lambda^4$, and
$C=\Lambda^2\oplus\Lambda^2\oplus\Lambda^2$. We use the matrices
\[
 L=\begin{pmatrix}
 1&0&1&0\\1&0&0&0\\0&1&0&1\\0&0&0&1\\
 2&1&1&1\\0&0&1&1
 \end{pmatrix},\qquad
 R=\begin{pmatrix}-1&-1&-1&0&1&0\\-1&1&0&-1&0&1\end{pmatrix}.
\]
Write $L_s$ for the two-row block of $L$ numbered $s=1,2,3$, and
$Q_s=L_s^{\mathsf t}L_s$. There is no finite three-colored graph $G$,
no positive integer $d$, and no direct-sum decomposition
\[
 H_1(G,\Lambda)=\iota(U)\oplus U'
\]
which is orthogonal for each of the three forms
$\sum_{e\text{ of color }s}x_e^2$ and whose restriction to $U$ is
$dQ_s$. Here $x_e$ is the oriented edge-coordinate function on graph
homology. Edge orientations do not affect the assertion.

\begin{proof}
We prove the assertion for every positive level. We adapt the map in
\citet[Proposition~6.14]{cyEngelFortmanSchreieder2026}.
We spell out the two-component version, including the map from their
Lemmas 6.8--6.13. We do not apply their regular-matroid theorem to
this nonregular matrix.

\emph{1. The exact profile identities.}
We have $RL=0$, and the first four rows of $L$ form a unimodular
matrix. Since the last two columns of $R$ form the identity,
\[
 0\longrightarrow\mathbb Z^4\xrightarrow{L}\mathbb Z^6
       \xrightarrow{R}\mathbb Z^2\longrightarrow0
 \tag{CG.1}
\]
is split exact. Each $L_s:\mathbb Z^4\to\mathbb Z^2$ is onto:
the first two blocks exhibit coordinate projections after unimodular
changes, and the last block has a minor of determinant one in
columns two and three. All these assertions remain true over
$\Lambda$.

Suppose a splitting as in the statement exists. We may join the
components of $G$ by edges in a tree; their homology coordinates
are zero, so this changes none of the data. Hence we assume $G$
connected. Let $q_e=x_e\circ\iota\in U^*$ for an edge of color $s$.
If $L_su=0$, positivity over $\mathbb R$ gives
\[
                 0=dQ_s(u,u)=\sum_{e\text{ of color }s}q_e(u)^2.
\]
Thus $q_e$ factors through $L_s$. Since $L_s$ is integrally onto,
there is a unique $v_e\in\Lambda^2$ with
\[
                 q_e=v_e^{\mathsf t}L_s.
\]
For $s=1,2,3$, define the two-component characteristic chain map
\[
 T_s:\Lambda^2\longrightarrow C_1(G,\Lambda),\qquad
 T_s(z)=\sum_{e\text{ of color }s}(v_e^{\mathsf t}z)e.
\]
We then have
\[
 \iota(u)=\sum_sT_s(L_su),\qquad
 \sum_{e\text{ of color }s}v_ev_e^{\mathsf t}=dI_2.
 \tag{CG.2}
\]
The second identity follows by restricting the graph form and using
the surjectivity of $L_s$.

Let $\beta:H_1(G,\Lambda)\to U$ be the projection of the supposed
splitting. Define the profile map $\Phi:C_1(G,\Lambda)\to C$ by
placing $v_e$ in block $s$ on an edge of color $s$. Orthogonality
implies, for every $h\in\ker\beta$,
\[
 (L_su)^{\mathsf t}\sum_{e\text{ of color }s}h_ev_e=0
                             \qquad(u\in U).
\]
Surjectivity of $L_s$ gives the vanishing of that block profile.
Together with (CG.2), this proves the exact identity
\[
                  \Phi(h)=dL\beta(h)
                            \quad(h\in H_1(G,\Lambda)).
 \tag{CG.3}
\]
This is where the common orthogonal retraction, rather than only
an embedding with prescribed Gram matrices, is used.

\emph{2. Axis refinement and block-valued chains.}
Write $d=2^jd'$ with $d'$ odd, and fix $r=j+2$.
Replace each edge $e$ of color $s$ by an axis path whose two
nonnegative integer coordinate lengths are congruent to the two
entries of $v_e$ modulo $2^r\Lambda$. We may take each length
strictly positive, replacing a zero residue by $2^r$.
Orient the path in the positive directions of its two block axes.
Each of the two chains $T_s(e_1),T_s(e_2)$ is refined by keeping
its original coefficient constant along this entire path.
The resulting graph is a subdivision of $G$.

For every $u\in U$ the sum of these refined chains with coefficients
$(L_su)_a$ is still closed. The unweighted axis profile of every
integral cycle lies in
\[
              2^jL\mathbb Z^4+2^r\mathbb Z^6.
 \tag{CG.4}
\]
Indeed its profile agrees with (CG.3) modulo $2^r$; intersection
with $\mathbb Z^6$ gives (CG.4), by (CG.1) and localization at~2.
Choosing a vertex as origin therefore maps the refined graph to
the unit-axis graph of the real torus
\[
 \mathcal T_{r,j}=\mathbb R^6/P_{r,j},\qquad
 P_{r,j}=2^jL\mathbb Z^4+2^r\mathbb Z^6.
 \tag{CG.5}
\]
Explicitly, a vertex is sent to the profile of a path from the
origin, modulo $P_{r,j}$. Condition (CG.4) makes this independent
of the path. Every axis edge is sent to the corresponding unit
axis edge.

After pushforward, denote the two chains supported on the two
axes of block $s$ by $b_{s,1},b_{s,2}$. They need not be supported
on their individually numbered axes. We only require block
support, and we retain the exact closure condition
\[
       \sum_{s=1}^3\sum_{a=1}^2(L_su)_a b_{s,a}
                         \text{ is closed for every }u\in U.
 \tag{CG.6}
\]
Let $M_s$ be the $2\times2$ matrix whose $a$th column is the
block-$s$ axis profile of $b_{s,a}$. Refinement and (CG.2) give
\[
                           M_s\equiv dI_2\pmod{2^r\Lambda}.
 \tag{CG.7}
\]
No symmetry or positivity of the refined profile matrices is
assumed.

\emph{3. Coordinate rounding preserves every block profile matrix.}
For one real coordinate, define a continuous function $h_j$ by
\[
 h_j(n2^j+t)=
 \begin{cases}
 n2^j+2^jt,&0\le t\le1,\\
 (n+1)2^j,&1\le t\le2^j,
 \end{cases}
 \qquad n\in\mathbb Z.
 \tag{CG.8}
\]
For $j=0$ this is the identity. Apply it coordinatewise on
$\mathbb R^6$. It is equivariant under $2^j\mathbb Z^6$ and
hence under $P_{r,j}$. The straight-line homotopy to the identity
is equivariant as well, so it induces a map $H$ of
$\mathcal T_{r,j}$ homotopic to the identity.

On the unit-axis graph, $H$ sends an edge either to a point or
to a chain of $2^j$ unit edges along the same axis. Its image is
the subdivision of the coarse-axis graph whose vertices are
$2^j\mathbb Z^6/P_{r,j}$ and whose edges have length $2^j$.
Distinct coarse axes do not meet at an interior subdivision
vertex: this follows by reducing coordinates modulo $2^j$, since
$P_{r,j}\subset2^j\mathbb Z^6$.

The pushed chains $H_*b_{s,a}$ retain block support and (CG.6).
Their block profile matrices are exactly $M_s$, not merely
congruent to them. To prove this, fix $s,a$ and choose $u\in U$
with $L_su=e_a$. The closed chain in (CG.6) has block-$s$
profile equal to the $a$th column of $M_s$. Homotopy to the
identity preserves its homology class in $\mathcal T_{r,j}$ and
hence its full displacement vector in
$P_{r,j}\otimes\Lambda\subset\Lambda^6$. Other blocks have
zero block-$s$ profile, before and after $H$. This isolates the
chosen column and proves the assertion.

At an interior subdivision vertex of a coarse edge in block $s$,
only its incoming and outgoing edges occur. Let $v^-,v^+\in
\Lambda^2$ be their two chain-coefficient vectors. Condition
(CG.6) says
\[
                   (L_su)^{\mathsf t}(v^--v^+)=0
                                    \qquad(u\in U).
\]
Thus $v^-=v^+$, again by surjectivity of $L_s$. We can collapse
each chain of $2^j$ edges to one coarse edge, keeping this
constant two-component coefficient. After scaling all coordinates
by $2^{-j}$ we obtain block-valued chains $b'_{s,a}$ on the
unit-axis graph of
\[
             \mathbb R^6/(L\mathbb Z^4+4\mathbb Z^6).
 \tag{CG.9}
\]
They still satisfy (CG.6). Their block profile matrices satisfy
\[
          M'_s=2^{-j}M_s\in\operatorname{Mat}_2(\Lambda),
          \qquad M'_s\equiv d'I_2\pmod{4\Lambda}.
 \tag{CG.10}
\]
The divisibility in the first assertion also follows directly
from the constant coefficients along each collapsed chain.
This completes the all-level reduction. It does not require
the flow on an axis edge to be proportional to that axis.

\emph{4. The relaxed axis certificate modulo four.}
By (CG.1), the vertices of (CG.9) are identified with
$(\mathbb Z/4)^2$ via $R$. An axis edge of block $s$ has profile
$e_a$ and carries an arbitrary flux vector $v=(v_1,v_2)\in
\Lambda^2$. Its contribution to the embedding flow is
$L_s^{\mathsf t}v$. Define the vertex function
\[
                  f(x,y)=x-x^2+y^2-2xy\pmod4.
 \tag{CG.11}
\]
Multiply the first coordinate of the vector-valued boundary
equation (CG.6) at each vertex by $f(x,y)$ and sum. The result
is zero modulo four. We compute its edge contributions.

The two columns of $R$ in block one are $(-1,-1)$ and
$(-1,1)$. Along either positive axis edge, the change of $f$
is $1$ modulo four. The first coordinate of
$L_1^{\mathsf t}v$ is $v_1+v_2$. Thus the total contribution
of block one is the sum of all four entries of $M'_1$.
The first coordinate of $L_2^{\mathsf t}v$ is zero, so block
two contributes zero.

The columns in block three are $(1,0)$ and $(0,1)$, and the
first coordinate of $L_3^{\mathsf t}v$ is $2v_1$. Along the
first axis, the change of $f$ is $-2(x+y)$ and gives zero
after multiplication by $2v_1$. Along the second axis, the
change is $1+2(y-x)$ and gives $2v_1$ modulo four. The total
block-three contribution is therefore $2(M'_3)_{21}$.
We have proved
\[
          0\equiv\sum_{a,b=1}^2(M'_1)_{ab}
                         +2(M'_3)_{21}
             \equiv2d'\pmod4,
\]
contrary to the oddness of $d'$. This proves the statement.
\end{proof}

\emph{Geometric scope still to be established.}
The matrices above arise in the proposed correlated monodromy
calculation for a type-$(1,2)$ surface and its dual. The present
proof is an abstract graph obstruction only.
\citet[Lemma~2.13]{cyEngelFortmanSchreieder2026} gives a common
orthogonal splitting when an appropriate strict semistable curve
family over a common polydisc already exists. Its Theorem~4.10
constructs the requisite common graph after alterations for
regular rank-one matroidal families, using Proposition~4.5 and
the admissible-resolution arguments of Section~3. We have not
established that geometric transfer for this correlated rank-two
cone. Consequently the abstract theorem alone does not prove
nonalgebraicity of a Poincare class or refute C.4 for an abelian
surface.
